\abstract{РЕФЕРАТ}

Объем работы равен \formbytotal{lastpage}{страниц}{е}{ам}{ам}. Работа содержит \formbytotal{figurecnt}{иллюстраци}{ю}{и}{й}, \formbytotal{tablecnt}{таблиц}{у}{ы}{}, \arabic{bibcount} библиографических источников и \formbytotal{числоПлакатов}{лист}{}{а}{ов} графического материала. Количество приложений – 2. Графический материал представлен в приложении А. Фрагменты исходного кода представлены в приложении Б.

Перечень ключевых слов: веб-приложение, Django, Python, SQLite, Quill, OpenRouter, ИИ-ассистент, управление авторским контентом, книга, арка, глава, автосохранение, импорт PDF, импорт DOCX, экспорт DOCX, REST API.

Объектом разработки является веб-приложение для управления авторским контентом с интеграцией интеллектуального ассистента.

Цель работы -- разработка веб-приложения с иерархической организацией текста, встроенным редактором с автосохранением, импортом черновиков из PDF/DOCX, экспортом глав в DOCX и взаимодействием с ИИ-ассистентом.

В основе системы лежит клиент-серверная архитектура на Django с базой данных SQLite и интерфейсом на HTML/CSS. На клиентской стороне интегрирован редактор Quill, обеспечивающий автосохранение и подсчёт слов. Импорт черновиков из PDF и DOCX выполняется библиотеками PyPDF2 и python-docx с возможностью выбора диапазона страниц, экспорт глав в DOCX -- через python-docx с полным сохранением форматирования. ИИ-ассистент подключён через REST API OpenRouter, поддерживает контекстный и общий режимы, а также автоматически переключается между бесплатными моделями при достижении лимитов.

Проведенное функциональное тестирование показало, что система работает корректно и полностью соответствует техническому заданию. Разработанное веб-приложение может быть использовано авторами для написания и структурирования книг, импорта черновиков, экспорта готовых глав и получения помощи от ИИ.


\selectlanguage{english}
\abstract{ABSTRACT}
  
The volume of work is \formbytotal{lastpage}{page}{}{s}{s}. The work contains \formbytotal{figurecnt}{illustration}{}{s}{s}, \formbytotal{tablecnt}{table}{}{s}{s}, \arabic{bibcount} bibliographic sources and \formbytotal{числоПлакатов}{sheet}{}{s}{s} of graphic material. The number of applications is 2. The graphic material is presented in annex A. The layout of the site, including the connection of components, is presented in annex B.

List of keywords: web-application, Django, Python, SQLite, Quill, OpenRouter, AI assistant, authoring content management, book, arch, chapter, auto-preservation, PDF import, DOCX import, DOCX export, REST API

The object of development is a web-based content management application with an intelligent assistant.

The goal of the work is to develop a web application with hierarchical text organization, built-in editor with car preservation, import drafts from PDF/DOCX, export chapters to DOCX and interaction with an AI assistant.

The system is based on a client-server architecture in Django with SQLite database and HTML/CSS interface. The Quill editor is integrated on the client side, ensuring car safety and word counting. Importing drafts from PDF and DOCX is done by PyPDF2 and python-docx libraries with the ability to select a page range, exporting chapters in DOCX via python-docx with full formatting. The AI assistant is connected via the OpenRouter REST API, supports contextual and sharing modes, and automatically switches between free models when limits are reached.

Functional testing has shown that the system works correctly and fully meets the technical requirements. The web application developed can be used by authors to write and structure books, import drafts, export finished chapters, and obtain help from AI.
\selectlanguage{russian}
